\input{preamble}
\usepackage{eurosym}


\newcommand{\tableaucroise}[4]{
\begin{tabular}{|c|c|c|c|}
	\cline{2-4}
	\multicolumn{1}{c|}{} & #1 \\ \hline
	#2 \\ \hline
	#3  \\ \hline
	#4  \\ \hline
\end{tabular}
}

\begin{document}

\pagestyle{fancy}
\fancyhead[L]{Terminale STMG}
\fancyhead[C]{\textbf{TD n°7 : variables aléatoires}}
\fancyhead[R]{\today}

\exe{}{
On considère une variable aléatoire $X$ dont la loi de probabilité est donnée dans le tableau ci-dessous :
\begin{center}
\begin{tabular}{ | c || c | c | c | c | c |}
\cline{1-6}
$x_i$ & 1 & 2 & 3 & 4 & 5 \\ \hline
$P(X=x_i)$ & 0,1 & 0,2 & 0,4 & 0,25 & 0,05 \\ \hline
\end{tabular}
\end{center}
\begin{enumerate}
\item Calculer $P(X=1) + P(X=2) + P(X=3) + P(X=4) + P(X=5)$.
\item Est-il possible que $X=6$ ? Expliquer.
\item Calculer $P(X \leq 3)$.
\item Calculer l'espérance de la variable aléatoire $X$.
\end{enumerate}
}{exe:base}{
\begin{enumerate}
\item $P(X=1) + P(X=2) + P(X=3) + P(X=4) + P(X=5) = 1$.
\item Non car alors, $P(X=6)=0$.
\item $P(X \leq 3) = P(X=1)+P(X=2)+P(X=3) = 0,7$.
\item Selon la formule du cours :
\begin{align*} 
E(X) & = 1 \times P(X=1) + 2 \times P(X=2) + 3 \times P(X=3) + 4 \times P(X=4) + 5 \times P(X=5) \\
& = 0,1 + 0,4 + 1,2 + 1 + 0,25 = 2,95
\end{align*}
\end{enumerate}
}

%Volé au Mattuuuuh
\exe{}{
	On lance un dé équilibré. Si le résultat est 1, 2 ou 3, on perd 5 euros. Si le résultat est 4 ou 5, on gagne 2 euros. Sinon, on gagne 10 euros. 
	Posons $X$ la variable aléatoire correspondant au gain	du joueur.
	\begin{enumerate}
		\item
		Quelles sont les valeurs prises par $X$ ?
		\item
		Déterminer la loi de $X$.
		\item
		Calculer $E(X)$. Le jeu est-il équitable ?
	\end{enumerate}
}{exe:rv1}{
	\begin{enumerate}
		\item
		$X \in \set{-5; 2; 10}$.
		\item
		\begin{center}
		\begin{tabular}{ | c || c | c | c |}
		\cline{1-4}
		$x_i$ & -5 & 2 & 10 \\ \hline
		$P(X=x_i)$ & 3/6 & 2/6 & 1/6 \\ \hline
		\end{tabular}
		\end{center}
		\item
		$E(X) = -5 \times \dfrac36 + 2 \times \dfrac26 + 10 \times \dfrac16 = -\dfrac16$, le jeu n'est pas équitable.
	\end{enumerate}
}

\exe{}{
Un ami vous propose de jouer à pile ou face, seulement vous savez que sa pièce est déséquilibrée : 
la probabilité d'obtenir pile est égale à $\dfrac49$. \\
Les régles fixées sont les suivantes : si la pièce tombe sur pile votre ami vous donne 12 \euro{} sinon vous lui donnez 10 \euro.
On note $X$ la variable aléatoire représentant le montant de votre gain sur un lancer.
\begin{enumerate}
\item Quelles valeurs peut prendre la variable aléatoire $X$ ?
\item Calculer l'espérance de $X$.
\item Est-ce intéressant pour vous de jouer ? Justifier.
\end{enumerate}
}{exe:bernou}{
	\begin{enumerate}
		\item
		$X \in \set{-10 ; 12}$.
		\item
		$E(X) = -10 \times \dfrac59 + 12 \times \dfrac49 = -\dfrac29$.
		\item
		Le jeu n'est pas équitable.
	\end{enumerate}
}

%Manuel 
\exe{}{
Deux amis, Jean-Siméon L'angosiet et Sophia, s'entraînent au football. Sophia tire quatre fois de suite.
On estime qu'elle marque un but avec une probabilité égale à 0,7 et que ses tirs sont indépendants les uns des autres.
On note $X$ la variable aléatoire comptant le nombre de buts marqués par Sophia.
\begin{enumerate}
\item Donner la loi suivie par la variable $X$ ainsi que ses paramètres.
\item Donner la probabilité d'un chemin conduisant à l'évènement $\set{X = 2}$.
\item Combien y-a-t-il de tels chemins ?
\item Calculer $P(X=2)$.
\end{enumerate}
}{exe:binom1}{
\begin{enumerate}
\item $X \approx \text{Binom}(4;0,7)$.
\item La probabilité d'un chemin est $0,7^2 \times 0,3^2 \approx 0,0441$.
\item Il y en a $\binom{4}{2}$ soit 6.
\item $P(X=2) = \binom{4}{2} 0,7^2 \times 0,3^2 \approx 6 \times 0,0441 \approx 0,2646$.
\end{enumerate}
}

\newpage

\exe{}{
Le triangle de Pascal est un outil permettant de calculer facilement les coefficients binomiaux $\binom{n}{k}$ 
en s'appuyant sur les propriétés (vue en cours) :
\[ (i)~:~\binom{n}{0} = \binom{n}{n} = 1 \et (ii)~:~\binom{n}{k} = \binom{n-1}{k-1}+\binom{n-1}{k} . \]
On organise le triangle sous cette forme :
\begin{center}
\begin{tabular}{>{$}l<{$}|*{7}{c}}
\multicolumn{1}{l}{$n$} &&&&&&&\\\cline{1-1} 
0 &1&&&&&&\\
1 &1&1&&&&&\\
2 &1&2&1&&&&\\
3 &1&3&3&1&&&\\
4 &1&3+1&3+3&3+1&1&& \\ 
5 & \dots & \dots & \dots & \dots & \dots & \dots & \\
6 & \dots & \dots & \dots & \dots & \dots & \dots & \dots \\ \hline
\multicolumn{1}{l}{} &0&1&2&3&4&5&6\\\cline{2-8}
\multicolumn{1}{l}{} &\multicolumn{7}{c}{$k$}
\end{tabular}
\end{center}

Pour obtenir une nouvelle ligne, on complète les deux colonnes aux extrémités par des 1 (c'est ce que dit la propriété (i)), 
puis on complète la colonne n°$k$ en sommant les valeurs des colonnes $k$ et $k-1$ de la ligne précédente 
(c'est ce que dit la propriété (ii)). \\

\begin{enumerate}
\item En vous appuyant sur cette procédure, compléter les lignes 5 et 6 du triangle de Pascal.
\item Sans faire plus de calcul, donner $\binom{6}{4}$.
\item Remarquez-vous des symétries dans les lignes de ce triangle ? Comment les expliquer ?
\end{enumerate}
}{exe:pascal}{

\begin{enumerate}
\item
\begin{center}
\begin{tabular}{>{$}l<{$}|*{7}{c}}
\multicolumn{1}{l}{$n$} &&&&&&&\\\cline{1-1} 
0 &1&&&&&&\\
1 &1&1&&&&&\\
2 &1&2&1&&&&\\
3 &1&3&3&1&&&\\
4 &1&4&6&4&1&& \\ 
5 & 1 & 5 & 10 & 10 & 5 & 1 & \\
6 & 1 & 6 & 15 & 20 & 15 & 6 & 1 \\ \hline
\multicolumn{1}{l}{} &0&1&2&3&4&5&6\\\cline{2-8}
\multicolumn{1}{l}{} &\multicolumn{7}{c}{$k$}
\end{tabular}
\end{center}
\item On lit dans le tableau $\binom{6}{4} = 15$.
\item Oui des symétries par rapport au coefficients centerales qui sont dues à la formule :
\[ \binom{n}{k} = \dfrac{n!}{k!(n-k)!} = \dfrac{n!}{(n-k)!k!} = \binom{n}{n-k} . \] 
\end{enumerate}
}

\exe{}{
On considère une variable aléatoire $X$ suivant une loi binomiale de paramètres $\left ( 6 ; 3/5 \right)$. Donner :
\begin{enumerate}
\item les valeurs possibles pour $X$ ;
\item la probabilité associée à chacune de ces valeurs ;
\item l'espérance de $X$.
\end{enumerate}
}{exe:binom2}{
\begin{enumerate}
\item $X$ prend les valeurs de 0 à 6.
\item On a 
\begin{align*}
P(X=0) & = \binom{6}{0} \left ( \dfrac35 \right)^0 \left ( \dfrac25 \right)^6 \approx 0,0041 \\
P(X=1) & = \binom{6}{1} \left ( \dfrac35 \right)^1 \left ( \dfrac25 \right)^5 \approx 0,037 \\
P(X=2) & = \binom{6}{2} \left ( \dfrac35 \right)^2 \left ( \dfrac25 \right)^4 \approx 0,14 \\
\dots \\
P(X=6) & = \binom{6}{6} \left ( \dfrac35 \right)^6 \left ( \dfrac25 \right)^0 \approx 0,047
\end{align*}
\item On calcul :
\[ E(X) = 0 \times P(X=0) + 1 \times P(X=1) + 2 \times P(X=2) + 3 \times P(X=3) + 4 \times P(X=4) + 5 \times P(X=5) + 6 \times P(X=6) \]
on trouve : $\dfrac{18}5 \approx 3,6$.
\end{enumerate}
}

% volé au même gars
\exe{}{
	On pose un pion sur la case inférieure gauche d'un échiquier $4\times4$.
	On déplace le pion de case en case afin d'atteindre la case supérieure droite.
	Le pion peut soit se déplacer d'une case vers le haut, soit d'une case vers la droite. Il est donc impossible de faire une boucle.
	En combien de façons différentes peut-il atteindre la case supérieure droite ?
	
}{exe:binom3}{
Un tel chemin comprend 4 mouvements vers la droite et 4 mouvements vers le haut. 
Il y a donc 8 mouvements au total et, parmi ces mouvements 4 vers le droite, soit $\binom{8}{4} = 70$ chemins possibles.
}

%%%%%%%%%%%%%%%%%%%%%%%%%%%%%
%
\newpage
\fancyhead[C]{\textbf{Solutions - TD5}}
\shipoutAnswer

\end{document}